\chapter{Horizonte demográfico}

\titular{El perímetro pensionario crece \textbf{3,9~\% real} y la población de 65 años y más
crece \textbf{4,29~\%}. \textbf{Por persona de esa edad, el gasto cae 0,4~\% real.} El mismo paquete que aumenta \textbf{el perímetro
pensionario} en pesos lo reduce \textbf{por persona de 65 años y más}, y la diferencia
entera es demografía.}

\begin{table}[htbp]\centering\small
\caption{Indicadores demográficos, 2026--2027}\label{cua:demo}
\begin{tabular}{lrrr}
\toprule
Indicador & 2026 & 2027 & Var. (\%) \\
\midrule
Población total & 134.407.258 & 135.391.662 & +0,73 \\
Población de 65 años y más & 12.201.321 & 12.724.222 & +4,29 \\
Población en edad escolar básica (3--14) & 25.421.421 & 25.179.850 & $-$0,95 \\
Razón de dependencia adulta & 13,44 & 13,90 & --- \\
Índice de envejecimiento & 38,84 & 40,91 & --- \\
Tasa global de fecundidad & 1,84 & 1,82 & --- \\
Esperanza de vida al nacer & 75,86 & 76,04 & --- \\
\bottomrule
\end{tabular}
\\[2pt]\begin{minipage}{0.92\textwidth}\scriptsize Fuente: CONAPO, Proyecciones de la Población de México 2020--2070, población a mitad de año, cobertura nacional, tier \texttt{externa\_demografica}. \textbf{Este cuadro no contiene ninguna cifra presupuestal}, por la regla de separación de la nota de método.\end{minipage}
\end{table}


\quecambio{La \textbf{población total crece 0,73~\%} y la de \textbf{65 años y más, 4,29~\%}:
casi seis veces más rápido. La \textbf{población en edad escolar básica cae 0,95~\%}, unas
241.571 personas menos. El \textbf{índice de envejecimiento pasa de 38,84 a 40,91} y la
\textbf{razón de dependencia adulta de 13,44 a 13,90}. La \textbf{tasa global de fecundidad
baja de 1,84 a 1,82} hijos por mujer.}

\porqueimporta{Las dos presiones van en sentidos contrarios y \textbf{el presupuesto responde
a las dos con el mismo signo}: aumenta el gasto pensionario, cuya población crece, y aumenta
el gasto educativo, cuya población se reduce. \textbf{El resultado por persona es opuesto en
cada caso}, y sólo se ve al dividir. Un documento que se quede en los montos verá que los dos
crecen ---y que en tasa el FONE crece más---; \textbf{sólo al dividir aparece que uno cae por
persona y el otro sube casi ocho puntos.}}

\section{Declaración de perímetro}

\marginal{la regla de separación}
\textbf{Los cuadros presupuestales y los per cápita van separados, y así están.} El cuadro
de indicadores demográficos no contiene ninguna cifra de pesos; el de razones por persona no
contiene ningún agregado presupuestal. Es la regla de la nota de método y no se relaja aquí.

La capa demográfica es \textbf{precargada}, con tier \texttt{externa\_demografica}, y sus
ocho series se descargaron antes de la entrega del paquete. \textbf{Ninguna cifra del paquete
se mezcla con un denominador externo en la misma línea de texto.}

\section{Cuerpo}

\begin{table}[htbp]\centering\footnotesize
\caption{Razones por persona: cada una con su denominador declarado (pesos de cada año)}\label{cua:percap}
\begin{tabular}{lrrr}
\toprule
Razón & 2026 & 2027 & Var. real (\%) \\
\midrule
Perímetro pensionario / población 65 y más & 139.670 & 144.665 & $-$0,4 \\
FONE / población en edad escolar básica & 21.494 & 24.102 & +7,8 \\
Ramos 12 y 56 / población total & 1.781 & 1.967 & +6,2 \\
\midrule \multicolumn{4}{l}{\emph{El mismo numerador con dos denominadores distintos, para que se vea la diferencia}} \\
FONE 2027 / población en edad escolar básica 2027 & --- & 24.102 & --- \\
FONE 2027 / matrícula de básica, ciclo 2025--2026 & --- & 26.564 & --- \\
\bottomrule
\end{tabular}
\\[2pt]\begin{minipage}{0.92\textwidth}\scriptsize \textbf{Cada renglón declara denominador, año y cobertura.} Población: CONAPO, a mitad de año, nacional. Matrícula: SEP, ciclo 2025--2026, modalidad escolarizada, nacional. \textbf{Matrícula no es población en edad escolar}: los dos denominadores difieren en 10,2~\% y mueven la cifra por persona en esa misma proporción. \textbf{El primer renglón es un indicador de presión demográfica, NO una pensión media}: el denominador incluye a quien no recibe y el numerador puede incluir pensiones de menores de 65 años.\end{minipage}
\end{table}


\subsection*{Tres advertencias sobre los denominadores, y ninguna es retórica}

\textbf{Primera: la razón de pensiones no es una pensión media.} El denominador
---población de 65 años y más--- incluye a quien no recibe pensión alguna, y el numerador
---el perímetro del Anexo 3--- puede incluir pensiones de personas menores de esa edad.
\textbf{Los 144.665 pesos del cuadro son un indicador de presión demográfica y nada más.}
Leerlos como el monto que recibe una persona pensionada sería un error de perímetro de los
que este proyecto documenta desde 2020.

\textbf{Segunda: matrícula no es población en edad escolar.} El mismo FONE de 2027 dividido
entre la población de 3 a 14 años da \textbf{24.102 pesos} y dividido entre la matrícula de
educación básica del ciclo 2025--2026 da \textbf{26.564}. \textbf{Los dos denominadores
difieren en 10,2~\%} y la cifra por persona se mueve en esa misma proporción. La primera
responde «cuánto se gasta por persona en edad de estudiar»; la segunda, «cuánto por persona
inscrita». \textbf{Son preguntas distintas y no se mezclan en un cuadro.}

\textbf{Tercera: la matrícula disponible no es la del ejercicio.} El ciclo más reciente que
publica la SEP es 2025--2026 y el presupuesto es de 2027. \textbf{La razón por alumno se
publica una sola vez, con su año declarado}, y no se usa para calcular variaciones.

\subsection*{Dónde el paquete trae horizonte demográfico y dónde no}

La instrucción de esta corrida pide verificar si el paquete invoca la demografía cuando ésta
presiona al alza y la omite cuando podría liberar espacio fiscal. \textbf{Con la fuente
disponible la respuesta es parcial}: el CGPE proyecta finanzas públicas a 2032 y publica un
escenario macroeconómico de mediano plazo, \textbf{pero no publica una proyección
demográfica propia ni asocia sus supuestos de gasto a la evolución de la población}. La
población en edad escolar básica lleva años cayendo y \textbf{el documento no lo menciona} al
justificar el gasto educativo.

\textbf{No se afirma que la omisión sea deliberada.} Se afirma que el paquete no publica el
vínculo, y que sin él la lectura por persona hay que construirla desde fuera, como hace este
capítulo.

\subsection*{Un contraste que la fuente misma obliga a declarar}

Las proyecciones de CONAPO y las estimaciones del Banco Mundial difieren de forma
sistemática: entre 2018 y 2025 la diferencia está entre \textbf{1,05 y 1,16~\%} de la
población, siempre en el mismo sentido. \textbf{Todas las razones por persona de este
capítulo se moverían alrededor de un punto porcentual} si se usara la otra fuente. Se usa
CONAPO por consistencia con el marco institucional mexicano, y la diferencia se declara aquí.

\section{Qué no sabemos}

\begin{list}{}{\setlength{\leftmargin}{1.4em}\setlength{\itemsep}{2pt}}
\item[$\bullet$] \textbf{Cuántas personas reciben cada prestación pensionaria.} Es un padrón,
no un presupuesto. El padrón de Bienestar está en la capa precargada, pero cruzarlo con el
perímetro del Anexo 3 exigiría el analítico por programa, que respondió 404.
\item[$\bullet$] \textbf{El gasto en salud por persona sin seguridad social.} El registro de
IMSS-Bienestar es \textbf{de cobertura parcial} ---sólo las entidades adheridas--- y es un
padrón de registro, no una medición de quién carece de afiliación: quien no se registra no
aparece. \textbf{Usarlo como denominador daría una cifra falsa}, y por eso no está en el
cuadro.
\item[$\bullet$] \textbf{El gasto por afiliado del IMSS.} La serie disponible es de
\textbf{asegurados}, que cuenta a la persona asegurada y no a sus familiares beneficiarios:
un gasto por afiliado con ese denominador daría una cifra varias veces mayor que la
verdadera. \textbf{Asegurados no es derechohabientes}, y por eso tampoco está en el cuadro.
\item[$\bullet$] \textbf{Las proyecciones de CELADE}, que no se pudieron obtener: la API de
CEPALSTAT responde por indicador pero su catálogo devuelve 404 y el identificador no se
descubrió sin adivinar. El contraste con el Banco Mundial es un \textbf{sustituto declarado}.
\item[$\bullet$] \textbf{La distribución por entidad federativa} de cualquiera de estas
razones. Los datos de población por entidad están en la capa; los presupuestales, no.
\end{list}

\begin{ficha}{horizonte demográfico}
\fila{Perímetro}{Indicadores nacionales de CONAPO; y tres razones por persona con numerador
presupuestal declarado: perímetro del Anexo 3, FONE del Anexo 22, y suma de los ramos 12 y
56 del Anexo 1. \textbf{Los cuadros presupuestales y los per cápita van en tablas separadas.}}
\fila{Línea de comparación}{Línea P en los tres numeradores presupuestales.}
\fila{Deflactor}{Del PIB, 1,040, del CGPE 2027, aplicado a las razones por persona. Véase el
recuadro de la nota de método.}
\fila{Año de los pesos}{Pesos corrientes de cada año en los niveles; pesos de 2027 en las
variaciones reales.}
\fila{PIB y añada}{No se usan razones a PIB en este capítulo.}
\fila{Denominadores}{Tres de CONAPO ---población total, de 65 años y más, y en edad escolar
básica---, de las Proyecciones 2020--2070, a mitad de año y con cobertura nacional. Uno de
la SEP: matrícula de educación básica, ciclo 2025--2026, modalidad escolarizada, nacional.
\textbf{Cada uno con su año y su cobertura, como exige la nota de método.}}
\fila{Fuentes}{De la capa demográfica precargada: los archivos de población nacional,
matrícula de la SEP y contraste de población, con el registro de fuentes de esa carpeta. Los
numeradores presupuestales, de la carpeta de datos de este documento.}
\fila{Qué no puede afirmarse con ellas}{Que una persona pensionada reciba 144.665 pesos al
año, ni que un alumno de básica cueste 26.564: son cocientes entre un agregado y una
población, no prestaciones ni costos unitarios. Tampoco que el gasto educativo por persona
crezca porque haya más recursos por alumno: crece en parte porque hay menos personas en edad
de estudiar.}
\end{ficha}
